\section*{Discussion}

The final experiments support a narrower and more defensible interpretation of complex layers than a simple
claim that voting should monotonically improve classification.
Across both datasets, the most useful information was typically concentrated in the first few neurons, but
the exact pattern depended on the base learner and on the feature-count regime induced by regularization.

The real-data benchmark showed that later neurons can be useful, but only modestly and not uniformly.
On \texttt{colon\_cancer}, \texttt{cpl} and \texttt{svm} reached their best mean accuracies after several neurons,
whereas \texttt{logreg} was already best at the first neuron in both feature-count regimes.
This suggests that sequential feature removal can expose complementary predictive structure for some learners,
but it should not be interpreted as a general guarantee that deeper voting improves accuracy.
The high-feature regime also showed that broader first neurons do not necessarily exhaust the constructible
sequence: although they selected more features initially, weaker effective regularization allowed later nonzero
neurons to be fitted in many splits.

The synthetic benchmark clarifies why these regimes differ.
When the first neuron is broad, as in high-feature \texttt{logreg}, it can recover a large fraction of the true
signal at once, but later neurons contribute little additional information.
When the first neuron is narrower, as in low-feature \texttt{cpl} or \texttt{svm}, recovery of informative
features is more gradual, and meaningful gains may emerge only after several neurons.
This makes neuron depth a useful analytical dimension for studying how a learner decomposes signal, rather
than merely an aggregation mechanism.
The low-feature \texttt{logreg} result also illustrates why later-depth summaries must be interpreted together
with coverage.
Its apparent improvement at $n=2$ reflects the subset of cross-validation splits in which a second neuron could
be fitted; within those splits, the unweighted two-neuron vote is identical to the first-neuron prediction under
the tie-breaking rule.
Thus, available-split counts are part of the result rather than only a plotting diagnostic.

Across both benchmarks, training-based weighted voting changed the conclusions little.
The weighted and uniform vote trajectories were usually close, indicating that the main effect came from the
successive sparse fits and the changing feature pool, not from the particular vote-weighting scheme.

The comparison also exposes a recurring trade-off between signal recovery and feature purity.
Broader neurons often recovered a larger share of informative variables, but they also selected many noisy
features.
Narrower neurons were cleaner, but they could require several layers before accumulating comparable signal.
This trade-off is central to the interpretation of later neurons and should be reported explicitly rather than
collapsed into a single accuracy score, in line with broader feature-selection literature
\cite{saeys2007review,nogueira2018stability}.

An additional practical observation is that realized depth itself is informative.
Some colon cancer configurations reached the nominal limit of 31 neurons, whereas the synthetic configurations
terminated after 3--13 neurons.
This suggests that the number of constructible neurons is not a trivial implementation detail; it reflects how
fast the learner exhausts separable structure under the imposed regularization.

Key limitations remain.
\begin{itemize}[noitemsep]
\item \texttt{colon\_cancer} is a small and classical microarray dataset, so its statistical variability is
high.
\item The synthetic benchmark is informative because it provides ground truth, but its conclusions remain
conditional on the chosen latent-source data-generation process.
\item Regularization values were calibrated to obtain comparable first-neuron widths, not to equalize the
effective penalty across models; direct numerical comparison of $C$ between learners is therefore not
meaningful.
\item The present study does not perform nested model selection, so the reported settings should be interpreted
as controlled analysis regimes rather than globally optimal hyperparameters.
\end{itemize}

Future work should extend this analysis to additional high-dimensional biological datasets, investigate the
stability of neuron-wise feature recovery across repeated resampling \cite{lukaszuk2024importance}, and study whether alternative stopping
rules can identify the depth at which complementary signal is exhausted and later neurons become primarily
noise-driven.
